% ============================================================
% CAPÍTULO 3: CARBOHIDRATOS
% ============================================================
% CONTENIDO BASADO EN: Unidad de Aprendizaje II, Tema "Carbohidratos"
% PROPÓSITO: Especificar características, clasificación, estructura, funcionalidad e importancia biotecnológica de los carbohidratos.
% ============================================================

\section{Características generales y clasificación de los carbohidratos}

% --- CONSEJO: Define los carbohidratos y explica por qué son importantes para la vida (fuente de energía primaria, estructura). ---

\resumebox{definicion}{Carbohidratos}{Los carbohidratos, también llamados glúcidos, hidratos de carbono o azúcares, son biomoléculas compuestas principalmente por carbono (C), hidrógeno (H) y oxígeno (O). Su fórmula general es $\ce{(CH2O)_n}$.}

\subsection{Clasificación de los carbohidratos}

% --- CONSEJO: Clasifica los carbohidratos según su tamaño (número de unidades de azúcar). ---
\begin{enumerate}
    \item \textbf{Monosacáridos:} Son los carbohidratos más simples. No pueden hidrolizarse a compuestos más sencillos.
    \item \textbf{Disacáridos:} Formados por la unión de dos monosacáridos mediante un enlace glucosídico.
    \item \textbf{Oligosacáridos:} Formados por la unión de 3 a 10 monosacáridos.
    \item \textbf{Polisacáridos:} Son polímeros formados por la unión de más de 10 unidades de monosacáridos.
\end{enumerate}

\section{Estructura y funcionalidad de los carbohidratos}

\subsection{Monosacáridos}

% --- CONSEJO: Describe la estructura básica de los monosacáridos (grupos funcionales: aldehído o cetona, y múltiples grupos hidroxilo). Menciona su clasificación como aldosas o cetosas, y su número de carbonos (triosas, pentosas, hexosas). ---

\subsubsection{Estructura de los monosacáridos}
% --- CONSEJO: Explica que los monosacáridos pueden existir en formas lineales y cíclicas (anillos de furanosa y piranosa) en solución. ---

\subsubsection{Función de los monosacáridos}
\begin{itemize}
    \item \textbf{Fuente de energía:} La glucosa es la principal fuente de energía para las células.
    \item \textbf{Estructural:} La ribosa y la desoxirribosa forman parte de los ácidos nucleicos (ARN y ADN).
    \item \textbf{Precursores:} Sirven como bloques de construcción para otros carbohidratos y biomoléculas.
\end{itemize}

% --- CONSEJO: Aquí podrías incluir una tabla con ejemplos de monosacáridos importantes (glucosa, fructosa, galactosa, ribosa, etc.) ---

\subsection{Disacáridos y Oligosacáridos}

% --- CONSEJO: Explica que se forman por la condensación de dos (o más) monosacáridos mediante un enlace glucosídico. ---
% --- Menciona ejemplos comunes como sacarosa (glucosa + fructosa), maltosa (glucosa + glucosa) y lactosa (glucosa + galactosa). ---

\subsection{Polisacáridos}

% --- CONSEJO: Clasifica los polisacáridos por su función (almacenamiento o estructural). ---
\begin{itemize}
    \item \textbf{Polisacáridos de almacenamiento:}
    \begin{itemize}
        \item \textbf{Almidón:} Polímero de glucosa en plantas (amilosa y amilopectina).
        \item \textbf{Glucógeno:} Polímero de glucosa en animales (principalmente en hígado y músculo).
    \end{itemize}
    \item \textbf{Polisacáridos estructurales:}
    \begin{itemize}
        \item \textbf{Celulosa:} Polímero de glucosa en plantas (principal componente de la pared celular).
        \item \textbf{Quitina:} Polímero de N-acetilglucosamina en el exoesqueleto de artrópodos y pared celular de hongos.
    \end{itemize}
\end{itemize}

\section{Importancia biotecnológica de los carbohidratos}

% --- CONSEJO: Esta es una sección clave para la asignatura. Relaciona los carbohidratos con la biotecnología. ---
\begin{itemize}
    \item \textbf{Producción de biocombustibles:} La hidrólisis de polisacáridos (celulosa, almidón) y la fermentación de azúcares simples produce etanol y otros biocombustibles.
    \item \textbf{Industria alimentaria:} Uso de carbohidratos como edulcorantes, espesantes y estabilizantes.
    \item \textbf{Farmacéutica:} Uso de oligosacáridos y polisacáridos como agentes terapéuticos o como excipientes en medicamentos.
    \item \textbf{Biomateriales:} La celulosa y la quitina se utilizan para producir materiales biodegradables y biocompatibles.
\end{itemize}

\section{Técnicas de identificación de carbohidratos}

% --- CONSEJO: El programa pide enlistar las técnicas. Menciona aquí las pruebas cualitativas y cuantitativas más comunes. ---
\begin{itemize}
    \item \textbf{Pruebas cualitativas:}
    \begin{itemize}
        \item \textbf{Reacción de Molisch:} Prueba general para carbohidratos.
        \item \textbf{Prueba de Benedict:} Detecta azúcares reductores (monosacáridos y algunos disacáridos).
        \item \textbf{Prueba de la yodo (I\textsubscript{2}):} Detecta polisacáridos como el almidón y el glucógeno (forman un complejo de color azul-negro o pardo-rojizo).
        \item \textbf{Prueba de Seliwanoff:} Distingue entre cetosas y aldosas.
    \end{itemize}
    \item \textbf{Técnicas cuantitativas:}
    \begin{itemize}
        \item \textbf{Cromatografía (TLC, HPLC):} Separa e identifica diferentes carbohidratos.
        \item \textbf{Espectrofotometría:} Mide la absorbancia de compuestos coloreados formados en reacciones específicas para cuantificar la concentración.
    \end{itemize}
\end{itemize}

% --- CONSEJO: Se sugiere una práctica de laboratorio para demostrar las propiedades físicas y químicas de los carbohidratos según su clasificación. ---